\section{Discussion and Possible Outcomes}

This section discusses the range of possible outcomes from the proposed study, their implications for the academic literature and for DeFi protocol design, and the boundary conditions that constrain the interpretation and generalizability of the results. Unlike studies that present a single hypothesized outcome, this discussion acknowledges that the results could take several forms, each with distinct implications.

\subsection{Scenario 1: Behavioral Process Variables Matter (H2a Supported)}

The most favorable outcome from the perspective of the study's theoretical contribution is that behavioral process variables contain statistically significant and economically meaningful incremental predictive power for liquidation events, above and beyond conventional on-chain risk indicators. If this scenario materializes, the results would indicate that \emph{how} a borrower's position arrived at its current state—the sequence of active adjustments, the timing of responses to risk signals, and the consistency of risk management behavior—carries information about future outcomes that is not fully captured by what the position's current state happens to be. This finding would be consistent with a growing body of evidence in financial economics suggesting that the process of portfolio management, not merely the resulting portfolio composition, contains information about investor sophistication and future performance.

The contribution of such a finding would be threefold. First, at the theoretical level, it would provide the first systematic evidence that prospect theory's predictions—particularly those concerning loss aversion and reference-point effects—are observable in DeFi lending behavior at a granular, position-by-position level. The unique features of the DeFi setting—objective, protocol-defined reference points; complete behavioral observability; and the absence of confounding factors associated with traditional financial intermediation—would make this evidence particularly compelling from a research design perspective and would contribute to the broader program of testing behavioral finance theories in naturalistic settings.

Second, at the empirical level, it would introduce a new category of predictor variables—behavioral process variables—to the on-chain credit risk toolkit. While the incremental predictive power of these variables might be modest in absolute terms (as is typically the case when building on a strong baseline model), even a small but statistically significant improvement in predictive power could be practically relevant in the context of DeFi protocols, where even marginal improvements in liquidation prediction can have substantial economic consequences given the volumes involved.

Third, at the practical level, the results would suggest that DeFi protocols could complement their existing, purely state-dependent risk assessment mechanisms with behavioral early-warning indicators. For example, a protocol might monitor the rate at which borrowers in a specific health factor band are actively adjusting their positions, treating an abnormally low rate of adjustment among risky positions as an aggregate signal of forthcoming liquidation pressure.

\subsection{Scenario 2: Behavior is Systematic but Not Predictive (RQ1 Supported; RQ2 Not Supported)}

An alternative possible outcome is that the behavioral hypotheses of RQ1 are supported—borrowers do exhibit systematic patterns of active adjustment as their positions approach risky thresholds, and these patterns are consistent with the predictions of prospect theory—but that RQ2 is not supported because behavioral process variables do not provide statistically significant incremental predictive power beyond conventional risk indicators. In this scenario, the findings would demonstrate that borrower behavior near the liquidation threshold is structured rather than random, and that this structure can be described using the theoretical framework of prospect theory, but that the information contained in this behavioral structure is already fully captured by simpler, more direct risk indicators. The health factor itself, possibly supplemented by its recent trajectory and the current market volatility environment, might be a sufficient statistic for the credit-relevant information contained in borrower behavior.

While this scenario would fall short of the study's more ambitious contribution, it would nonetheless represent a valuable addition to the literature. It would provide what is, to the best of our knowledge, the first systematic empirical characterization of how DeFi borrowers actively manage—or fail to manage—their positions as risk accumulates. This empirical contribution would be valuable even in the absence of predictive success, as it would document a set of stylized facts about borrower behavior that could inform future theoretical and empirical work.

\subsection{Scenario 3: No Systematic Behavior or Predictive Power (RQ1 and RQ2 Not Supported)}

The least favorable outcome from the perspective of publication prospects is that neither RQ1 nor RQ2 is supported: borrowers do not exhibit systematic patterns of active adjustment as their positions approach the liquidation threshold, and behavioral process variables do not improve predictions of liquidation events. In this scenario, the data would suggest that DeFi borrowers' behavior near the liquidation threshold is either random, driven entirely by the mechanical effects of price movements and liquidation bot activity, or dominated by factors that the study is unable to observe.

Even in this scenario, the study would have produced a useful null result. The transparency of DeFi protocols means that the research design and data are fully replicable, and the hypotheses are falsifiable and were specified ex ante. However, a null result in this setting is more likely to reflect genuine absence of an effect than other frequently-invoked explanations for non-replication in the social sciences, such as researcher degrees of freedom or publication bias. Indeed, one plausible interpretation of this outcome is that DeFi lending markets are relatively efficient, with borrowers responding to the strong financial incentives created by liquidation penalties in a manner that is well-approximated by standard risk measures. This interpretation would itself be a finding of interest to researchers studying the efficiency of DeFi markets relative to traditional financial markets, where behavioral biases are more consistently documented.

\subsection{Contributions and Implications}

The proposed study has the potential to make contributions at multiple levels. At the methodological level, it would demonstrate a template for how publicly available blockchain data can be used to conduct detailed, process-oriented behavioral research without relying on proprietary datasets or platform partnerships. At the theoretical level, it would provide a systematic test of behavioral finance predictions in a novel empirical setting that offers unusually clean identification of the relevant theoretical constructs. At the practical level, it would inform the design of DeFi protocols by evaluating whether borrower behavior contains credit-relevant information that could complement existing, state-based risk assessment mechanisms.

The implications of the study extend beyond the specific context of DeFi lending. The broader question addressed—whether the \emph{process} of risk management contains information beyond the \emph{state} of risk exposure—applies to a wide range of financial contexts, from margin trading in traditional securities markets to the management of corporate debt structures. Demonstrating that this question can be answered in the affirmative using DeFi data would provide a proof of concept for a research program that could eventually extend to these other contexts, where the necessary data would be considerably more difficult to obtain.

\subsection{Limitations and Boundary Conditions}

The proposed study has several important limitations that constrain the interpretation and generalizability of its results, regardless of which of the three scenarios described above materializes.

\textbf{Partial observability of user behavior.} The analysis is limited to on-chain transactions, which means that it cannot observe borrower behavior on centralized exchanges or off-chain activities. A borrower who appears to be relatively passive on-chain may be actively hedging their position on a centralized exchange, and a borrower who appears to be actively adjusting their position on-chain may be doing so for reasons unrelated to the risk of that specific position. This partial observability is an inherent limitation of research using public blockchain data and should be explicitly acknowledged rather than minimized.

\textbf{Cannot infer intent from on-chain data.} While the classification of transactions as active borrower actions or passive liquidations is straightforward and replicable, the motives underlying those actions cannot be observed. A borrower who adds collateral following a decline in their health factor may be responding to the risk of liquidation, but they could also be responding to factors that are correlated with—but conceptually distinct from—position risk, such as a change in their assessment of the collateral asset's future price trajectory or a desire to maintain a specific portfolio allocation.

\textbf{Limited external validity.} The study is focused on a specific set of DeFi lending protocols during a specific period of market development. The extent to which the findings would generalize to other protocols, other time periods, or other types of blockchain-based financial applications is unknown and would need to be evaluated through replication and extension. The rapid pace of change in the DeFi ecosystem—including ongoing changes to protocol parameters, the introduction of new lending mechanisms, and the evolution of the competitive landscape—means that even findings that are internally valid at the time of analysis may have a limited shelf life.

\textbf{Confounding by unobserved borrower characteristics.} While panel data methods with borrower fixed effects will be used to control for time-invariant unobserved heterogeneity, it remains possible that time-varying unobserved factors—such as changes in a borrower's financial circumstances, risk preferences, or access to alternative sources of liquidity—could influence both their behavioral responses to risk accumulation and their subsequent likelihood of liquidation, creating a spurious association between behavioral process variables and liquidation outcomes. The observational nature of the data means that the study can identify associations but cannot establish causal relationships.

\subsection{Ethical Considerations}

The proposed research uses only publicly available on-chain data that does not contain personally identifiable information. No individual user will be contacted or re-identified from their wallet address, and the analysis will be conducted exclusively at the level of aggregates and statistical summaries. Nevertheless, the study touches on broader ethical questions about the use of public blockchain data for credit assessment. Following the work of this study, protocols might seek to incorporate behavioral variables into risk assessment, and this could have consequences for borrowers that are not easily anticipated. It is worth noting that the relationship between active position adjustment and credit outcomes is assessed at the aggregate level; the estimates cannot be meaningfully interpreted as individual-level predictions without careful attention to issues of statistical discrimination and fairness.

\subsection{Future Research Directions}

The proposed study opens several avenues for future research, some of which can be anticipated at this stage and others that will become apparent only after the empirical results are available. First, if the behavioral process variables are found to have incremental predictive power, subsequent research could investigate the mechanisms underlying this relationship using more detailed decomposition analyses or structural estimation. Second, the framework developed in this study could be extended to other DeFi lending protocols operating on different blockchain networks, providing evidence on the generalizability of the results across different protocol designs and market environments. Third, research could investigate the welfare implications of incorporating behavioral process variables into protocol risk parameters, considering both the potential benefits of improved risk assessment and the potential costs associated with increased protocol complexity. Fourth, future work could examine whether borrowers who exhibit systematic patterns of active position management in DeFi have different characteristics or outcomes in traditional credit markets, conditional on the availability of data that links the two domains.
